\chapter*{练习}
\addcontentsline{toc}{section}{练习}

\section*{练习 4.1}
\addcontentsline{toc}{subsection}{练习 4.1}

对于矩阵

\begin{equation*}
\boldsymbol{A} = \begin{bmatrix} 1 & 3 & 5 \\ 2 & 4 & 6 \\ 0 & 2 & 4 \end{bmatrix}
\end{equation*}

分别使用拉普拉斯展开（沿第一行）和萨鲁斯法则（Sarrus' rule）计算其行列式。

\medskip\hrule\medskip

\section*{练习 4.2}
\addcontentsline{toc}{subsection}{练习 4.2}

高效计算下列矩阵的行列式：

\begin{equation*}
\begin{adjustbox}{max width=\linewidth}
\begin{bmatrix} 2 & 0 & 1 & 2 & 0 \\ 2 & -1 & 0 & 1 & 1 \\ 0 & 1 & 2 & 1 & 2 \\ -2 & 0 & 2 & -1 & 2 \\ 2 & 0 & 0 & 1 & 1 \end{bmatrix}
\end{adjustbox}
\end{equation*}

\medskip\hrule\medskip

\section*{练习 4.3}
\addcontentsline{toc}{subsection}{练习 4.3}

计算下列矩阵的特征空间：

\begin{equation*}
\begin{bmatrix} 1 & 0 \\ 1 & 1 \end{bmatrix}, \quad \begin{bmatrix} -2 & 2 \\ 2 & 1 \end{bmatrix}
\end{equation*}

\medskip\hrule\medskip

\section*{练习 4.4}
\addcontentsline{toc}{subsection}{练习 4.4}

计算矩阵

\begin{equation*}
\begin{adjustbox}{max width=\linewidth}
\boldsymbol{A} = \begin{bmatrix} 0 & -1 & 1 & 1 \\ -1 & 1 & -2 & 3 \\ 2 & -1 & 0 & 0 \\ 1 & -1 & 1 & 0 \end{bmatrix}
\end{adjustbox}
\end{equation*}

的所有特征空间。

\medskip\hrule\medskip

\section*{练习 4.5}
\addcontentsline{toc}{subsection}{练习 4.5}

矩阵的可对角化性与其可逆性无关。判断下列四个矩阵是否可对角化和/或可逆：

\begin{equation*}
\begin{bmatrix} 1 & 0 \\ 0 & 1 \end{bmatrix}, \quad \begin{bmatrix} 1 & 0 \\ 0 & 0 \end{bmatrix}, \quad \begin{bmatrix} 1 & 1 \\ 0 & 1 \end{bmatrix}, \quad \begin{bmatrix} 0 & 1 \\ 0 & 0 \end{bmatrix}
\end{equation*}

\medskip\hrule\medskip

\section*{练习 4.6}
\addcontentsline{toc}{subsection}{练习 4.6}

计算下列变换矩阵的特征空间。它们可对角化吗？

a.

\begin{equation*}
\boldsymbol{A} = \begin{bmatrix} 2 & 3 & 0 \\ 1 & 4 & 3 \\ 0 & 0 & 1 \end{bmatrix}
\end{equation*}

b.

\begin{equation*}
\begin{adjustbox}{max width=\linewidth}
\boldsymbol{A} = \begin{bmatrix} 1 & 1 & 0 & 0 \\ 0 & 0 & 0 & 0 \\ 0 & 0 & 0 & 0 \\ 0 & 0 & 0 & 0 \end{bmatrix}
\end{adjustbox}
\end{equation*}

\medskip\hrule\medskip

\section*{练习 4.7}
\addcontentsline{toc}{subsection}{练习 4.7}

下列矩阵是否可对角化？若是，求其对角形式以及变换矩阵关于其呈对角形式的基。若否，给出它们不可对角化的理由。

a.

\begin{equation*}
\boldsymbol{A} = \begin{bmatrix} 0 & 1 \\ -8 & 4 \end{bmatrix}
\end{equation*}

b.

\begin{equation*}
\boldsymbol{A} = \begin{bmatrix} 1 & 1 & 1 \\ 1 & 1 & 1 \\ 1 & 1 & 1 \end{bmatrix}
\end{equation*}

c.

\begin{equation*}
\begin{adjustbox}{max width=\linewidth}
\boldsymbol{A} = \begin{bmatrix} 5 & 4 & 2 & 1 \\ 0 & 1 & -1 & -1 \\ -1 & -1 & 3 & 0 \\ 1 & 1 & -1 & 2 \end{bmatrix}
\end{adjustbox}
\end{equation*}

d.

\begin{equation*}
\boldsymbol{A} = \begin{bmatrix} 5 & -6 & -6 \\ -1 & 4 & 2 \\ 3 & -6 & -4 \end{bmatrix}
\end{equation*}

\medskip\hrule\medskip

\section*{练习 4.8}
\addcontentsline{toc}{subsection}{练习 4.8}

求矩阵

\begin{equation*}
\boldsymbol{A} = \begin{bmatrix} 3 & 2 & 2 \\ 2 & 3 & -2 \end{bmatrix}
\end{equation*}

的奇异值分解（SVD）。

\medskip\hrule\medskip

\section*{练习 4.9}
\addcontentsline{toc}{subsection}{练习 4.9}

求矩阵

\begin{equation*}
\boldsymbol{A} = \begin{bmatrix} 2 & 2 \\ -1 & 1 \end{bmatrix}
\end{equation*}

的奇异值分解。

\medskip\hrule\medskip

\section*{练习 4.10}
\addcontentsline{toc}{subsection}{练习 4.10}

求矩阵

\begin{equation*}
\boldsymbol{A} = \begin{bmatrix} 3 & 2 & 2 \\ 2 & 3 & -2 \end{bmatrix}
\end{equation*}

的最佳秩 1 近似。

\medskip\hrule\medskip

\section*{练习 4.11}
\addcontentsline{toc}{subsection}{练习 4.11}

证明：对于任意 $ \boldsymbol{A} \in \mathbb{R}^{m \times n} $，矩阵 $ \boldsymbol{A}^\top \boldsymbol{A} $ 与 $ \boldsymbol{A} \boldsymbol{A}^\top $ 具有相同的非零特征值。

\medskip\hrule\medskip

\section*{练习 4.12}
\addcontentsline{toc}{subsection}{练习 4.12}

证明对于 $ \boldsymbol{x} \neq \boldsymbol{0} $ 定理 4.24 成立，即证明：

\begin{equation*}
\max_{\boldsymbol{x}} \frac{\|\boldsymbol{A} \boldsymbol{x}\|_2}{\|\boldsymbol{x}\|_2} = \sigma_1
\end{equation*}

其中 $ \sigma_1 $ 是 $ \boldsymbol{A} \in \mathbb{R}^{m \times n} $ 的最大奇异值。
